\chapter{REVIEW OF RELATED LITERATURE}\label{chap:review-of-related-literature}

This chapter reviews the literature that supports the development of the
proposed TinyML Smart Plug as a socket-level protection device for
progressive electrical-fire precursors. Rather than treating electrical
faults as a single undifferentiated class of events, the review is
organized around the specific hazards that motivated the prototype:
overload-related thermal stress, localized contact heating at the
plug-socket interface, series arc faults, and abnormal supply conditions
that may intensify equipment stress or affect protection behavior. It
also reviews related literature on non-intrusive load monitoring, TinyML
deployment on microcontrollers, current-signal feature extraction,
embedded fault detection, and validation practices for
protection-oriented devices. The chapter is divided into three major
parts: theoretical background, related studies, and synthesis.

\section{Theoretical Background}\label{theoretical-background}

This section presents the theoretical concepts needed to understand the
fault conditions addressed by the study. It begins with the residential
electrical-fire problem and then narrows the discussion to electrical
fault mechanisms, conventional protective devices, contact heating,
external temperature measurement, waveform disturbance, conduction
intermittence, and load-family behavior. These concepts establish why
the proposed system was designed as a complementary socket-level
protection device rather than as a replacement for existing circuit
protection.

\subsection{Residential Fire-Risk and Protection Gap}\label{residential-fire-risk-and-protection-gap}

Electrical faults remain a persistent residential safety problem because
not all hazardous conditions develop as instant catastrophic failures.
In the Philippines, {[}1{]} identified electrical issues as the leading
cause of recorded fire incidents, emphasizing that domestic electrical
safety remains a practical and not merely theoretical concern. This is
reinforced by NFPA fire statistics showing that electrical failure or
malfunction continues to account for a substantial share of home
structure fires, injuries, and losses {[}2{]}.

The core implication for this study is that many dangerous electrical
events evolve gradually. Some are cleared quickly by conventional
protection, but others develop in a sub-breaker regime in which heat,
contact degradation, and intermittent conduction are already becoming
hazardous while the branch circuit still appears electrically
acceptable. This gap between ordinary operation and fully developed
failure is the operating space addressed by the proposed smart plug.

Since residential fire risk can arise from different electrical
mechanisms, the next section distinguishes among the fault classes that
may occur in household circuits and clarifies which of these are most
relevant to the proposed device.

\subsection{Electrical Fault Mechanisms}\label{electrical-fault-mechanisms}

Overload is a sustained current demand above the intended operating
capacity of a conductor, outlet, or protection device. It is
fundamentally a thermal process because Joule heating rises with current
and time. Conventional low-voltage breakers are intentionally
inverse-time devices: lower overcurrents are tolerated longer, while
larger overcurrents trip more quickly {[}12{]}. This design is necessary
to avoid nuisance trips, but it also means that a localized plug, relay,
receptacle, or appliance cord can remain under significant thermal
stress before upstream protection is forced to operate. For common
PVC-insulated low-voltage cables, the 70°C conductor temperature remains
a meaningful thermal reference, and repeated thermal stress accelerates
insulation aging and deterioration {[}13{]}, {[}14{]}.

Short circuits and ground faults are more abrupt fault classes. A short
circuit produces a low-impedance path and very high current, while a
ground fault diverts current through an unintended path to earth or
grounded metal. For personnel protection, the danger of ground-fault
current is severe because fibrillation can occur at relatively low
current levels, especially when body resistance is reduced by moisture
{[}15{]}. In practice, these fault classes are already well served by
breakers and residual-current or ground-fault protective devices, so
they are reviewed here mainly to clarify why they were not the primary
target of the present smart plug.

Arc faults occupy a more difficult middle ground. A parallel arc can
resemble a short circuit and may trip conventional overcurrent
protection, but a series arc is limited by the load current and may
therefore remain below the normal breaker trip region. AFCI literature
consistently identifies this as the main reason arc-fault protection was
developed: the waveform can be thermally dangerous without becoming a
classical overcurrent event {[}16{]}. For a smart plug intended to
protect the immediate load connection, this distinction is crucial
because many plug-side and cord-side failures are series-arc-like before
they ever become hard faults.

High-resistance connection faults are especially relevant to this study
because they explain why socket-level temperature monitoring is needed.
A loose, corroded, worn, oxidized, or mechanically weakened contact
increases local resistance, so the same operating current generates
disproportionate local heating according to P=I\textsuperscript{2}R.
Classic NIST work on glowing electrical connections and later studies on
high-resistance mechanisms show that such faults can persist for long
periods, ignite nearby polymers, and remain effectively invisible to
ordinary breakers or GFCIs because the total current may stay within an
otherwise normal appliance range {[}17{]}, {[}18{]}.

Overvoltage and undervoltage are not the main fire mechanisms targeted
by the research prototype, yet they remain important because they
influence equipment stress and protection logic. Transient overvoltages
can damage connected electronics, which is why surge protective devices
are used to limit overvoltage at the load side {[}19{]}, {[}20{]}.
Sustained undervoltage can also be harmful because some loads,
especially motors, may overheat or operate inefficiently under depressed
supply conditions. In the context of this study, voltage supervision
therefore serves two purposes: direct protection against abnormal supply
conditions and better interpretation of current-based fault signatures.

After identifying the major electrical fault mechanisms, it is necessary
to examine how existing protective devices respond to these conditions.
This establishes why some faults are already adequately handled by
conventional protection, while others require complementary monitoring.

\subsection{Conventional Protective Devices and Their Limits}\label{conventional-protective-devices-and-their-limits}

Conventional protection is effective because it is selective, but that
same selectivity creates blind spots. Miniature circuit breakers
primarily protect branch wiring against overload and short circuit.
Their thermal-magnetic behavior, including fast tripping at high
multiples of rated current and delayed tripping at lower overcurrent
levels, is intentionally coordinated around conductor protection rather
than local receptacle or plug-blade temperature {[}3{]}, {[}21{]},
{[}22{]}.

Ground-fault circuit interrupters compare outgoing and returning current
and trip on small imbalances, typically around the milliampere range
associated with leakage or shock hazards {[}23{]}. AFCIs were introduced
to detect wiring arcs that overcurrent protection can miss, especially
series arcs in branch circuits {[}16{]}. Surge protective devices,
meanwhile, are intended to clamp transient overvoltages and prevent
damage to downstream equipment {[}19{]}, {[}20{]}.

Taken together, these devices provide important layers of protection,
but none directly measure the thermal condition of the plug-socket
interface. They are therefore not designed to identify a local hotspot
created by poor contact pressure, oxidation, or partial separation while
the total current remains modest. This is the central theoretical
justification for the proposed prototype: not replacement protection,
but complementary protection at the socket level, where early-stage
heating and unstable conduction can begin.

Because conventional protection does not directly observe the
plug-socket interface, the next section focuses on the thermal behavior
of poor contacts and degraded sockets, which form one of the main fault
targets of the study.

\subsection{Thermal Behavior of Poor Contacts and Degraded Sockets}\label{thermal-behavior-of-poor-contacts-and-degraded-sockets}

The literature on glowing contacts, poor plug-receptacle interfaces, and
overheated outlets converges on a single physical principle:
electrical-fire precursors are often thermally local before becoming
electrically global. The Development and Analysis of Electrical
Receptacle Fires, authored by {[}24{]}, documented that loose terminal
conditions resulted in substantial overheating, visible damage, and
failure in receptacle systems. Later fire studies showed that glowing
connections can persist, consume surrounding non-metal parts, and ignite
plastics even though conventional protection may not operate {[}7{]},
{[}17{]}.

The importance of a socket-level monitor is therefore not merely that it
measures temperature, but that it measures temperature near the region
where hidden degradation manifests first. The thermal process may be
modeled as a first-order system in which the local node temperature
rises progressively toward a steady-state value set by injected Joule
heating, thermal resistance to ambient, and thermal capacitance. In this
view, dangerous faults do not need to create instant flame to be
serious; they only need enough sustained local power to initiate
oxidation, polymer softening, charring, and further growth of contact
resistance.

Standards and fire studies also show that external body temperature is
not equivalent to true contact temperature. IEC 60884-1 classifies
household plugs and socket-outlets for ordinary ambient temperatures not
normally exceeding 40°C and uses temperature-rise criteria at internal
conductive points rather than prescribing one universal external body
limit {[}25{]}. IEC 60227 similarly preserves 70°C as a meaningful
conductor-temperature reference for common PVC-insulated cable families,
since this region is associated with insulation stress and visible wire
degradation {[}14{]}. UL 498 also provides temperature-rise limits for
attachment plugs and receptacles under test conditions {[}26{]}. When
these standards are interpreted together with thermal lag and hidden
contact-spot behavior, an early external intervention threshold becomes
more defensible for a non-intrusive prototype.

\subsection{External Temperature Measurement and Hotspot Underestimation}\label{external-temperature-measurement-and-hotspot-underestimation}

For embedded socket monitoring, the NTC thermistor is useful because it
is low-cost, sensitive in the temperature range of interest, easy to
interface with, and suitable for continuous protection duty {[}27{]},
{[}28{]}. However, the same literature also warns that NTC measurements
must be interpreted with care. {[}27{]} defines thermal time constant as
the time required for the sensor to traverse 63.2\% of a sudden
temperature change, while {[}28{]} notes that the thermistor is useful
for over-temperature protection but does not directly measure the
hottest internal point of the system. {[}29{]} further explains that
poor sensor placement increases lag because the controller only reacts
to the temperature it ``sees'' through the measurement point.

This matters directly for receptacle monitoring because the hottest
point is usually inside the contact interface, and not at the external
surface touched by the thermistor. {[}30{]} argued that surface
temperature is a poor severity indicator when overheated contact spots
are thermally separated from external surfaces. The study by {[}31{]} on
indirect thermography indicates that the internal temperature of
electromechanical relay contacts can significantly exceed the package
temperature. In other words, an external sensor reading near 40°C can
correspond to a hidden contact zone already much hotter than the
measured point.

For that reason, the 40°C threshold used in this study is best
interpreted as a conservative early intervention threshold for a
non-intrusive, thermally lagged sensor used to estimate a hidden
hotspot. It sits below the internal temperature region associated with
clearly abnormal accessible heating and well below severe
material-stress levels, while remaining sufficiently above normal
ambient temperatures. This provides a safety margin for sensor lag,
thermal gradients, and placement uncertainty.

With the thermal aspects of the fault problem established, we now turn
our attention to series arc faults. Unlike contact heating, which is
mainly thermal and localized, series arcs need to be detected based on
their waveform behavior, as they can occur at current levels that are
limited by the load.

\subsection{Waveform Disturbance and Conduction Intermittence}\label{waveform-disturbance-and-conduction-intermittence}

Arc fault detection is possible only because arcing modifies the current
waveform in ways that healthy operation does not. The literature
repeatedly describes arc currents as random, discontinuous, pulse-rich,
and rich in abnormal harmonic or high-frequency energy, but it also
warns that the exact appearance depends strongly on the connected load
{[}32{]}--{[}34{]}. This has two important implications for the study.

First, it justifies a waveform-disturbance family of features. These
include descriptors related to harmonic distortion, spectral change,
residual impulsiveness, and band-energy redistribution. In the
literature, THD, high-frequency noise, residual or singular components,
and spectral fluctuation are regularly used because arcs inject
non-stationary energy that does not behave like ordinary steady-state
appliance current {[}34{]}--{[}39{]}.

Second, it justifies a conduction-intermittence family of features.
Series arcs and unstable contacts often create partial conduction,
zero-current gaps, restrikes, half-cycle asymmetry, and irregular
zero-cross behavior. Accordingly, the literature emphasizes zero-current
features, pulse counting, current-fluctuation descriptors, and
zero-cross interval instability because they directly represent the
intermittent-recontact nature of arcing and poor electrical connection
{[}34{]}, {[}36{]}, {[}37{]}.

Viewed together, waveform disturbance and conduction intermittence form
a multi-domain profile of an arc fault. However, these indicators cannot
be interpreted in isolation because many normal household loads are also
electrically nonlinear. The next section therefore discusses the
importance of non-intrusive load monitoring and load-family context in
reducing false detections.

\subsection{Non-Intrusive Load Monitoring and Load Families}\label{non-intrusive-load-monitoring-and-load-families}

A further theoretical issue is that many healthy appliances are already
nonlinear. Motors produce inrush and commutation effects; rectifier and
SMPS loads draw peaky current; and phase-angle-controlled loads
intentionally chop the waveform. NILM research demonstrates that these
differences are not just nuisance artifacts but informative load
signatures that can be extracted from current and voltage waveforms for
appliance identification and load disaggregation {[}8{]},
{[}40{]}--{[}42{]}.

This is precisely why load family matters in fault detection. A waveform
that looks abnormal for a resistive heater may be normal for a universal
motor or rectifier load. Consequently, the detector must distinguish
fault-induced non-stationarity from load-intrinsic nonlinearity.
Context-aware inference, whether framed as NILM, load-family
recognition, or appliance-family conditioning, is therefore not optional
but a false-positive control strategy.

TinyML makes this context-aware approach feasible on embedded hardware.
Reviews of machine learning for microcontroller-class hardware and
TensorFlow Lite Micro show that lightweight models can run within tight
power, memory, and latency budgets when the feature space and runtime
are carefully designed {[}32{]}, {[}43{]}. This makes a hybrid embedded
architecture realistic: deterministic thresholds for hard safety limits,
with lightweight inference for ambiguous waveform cases such as
load-dependent series arcs.

To summarize the theoretical boundary of the study, Table 2.1 compares
the main electrical fault classes reviewed in this section. The table
identifies the primary hazard associated with each fault, whether
conventional protection generally handles it, and how each fault class
relates to the proposed smart plug.

\subparagraph{\texorpdfstring{\textbf{Table 2.1.} Summary of Electrical
Faults}{Table 2.1. Summary of Electrical Faults}}\label{table-2.1.-summary-of-electrical-faults}

{\def\LTcaptype{none} % do not increment counter
\begin{longtable}[]{@{}
  >{\raggedright\arraybackslash}p{(\linewidth - 6\tabcolsep) * \real{0.1994}}
  >{\raggedright\arraybackslash}p{(\linewidth - 6\tabcolsep) * \real{0.2685}}
  >{\raggedright\arraybackslash}p{(\linewidth - 6\tabcolsep) * \real{0.1849}}
  >{\raggedright\arraybackslash}p{(\linewidth - 6\tabcolsep) * \real{0.2781}}@{}}
\toprule\noalign{}
\begin{minipage}[b]{\linewidth}\raggedright
\textbf{Fault class}
\end{minipage} & \begin{minipage}[b]{\linewidth}\raggedright
\textbf{Primary hazard}
\end{minipage} & \begin{minipage}[b]{\linewidth}\raggedright
\textbf{Handled by Conventional Protection}
\end{minipage} & \begin{minipage}[b]{\linewidth}\raggedright
\textbf{Relevance}
\end{minipage} \\
\begin{minipage}[b]{\linewidth}\raggedright
Overload
\end{minipage} & \begin{minipage}[b]{\linewidth}\raggedright
Thermal stress, insulation degradation, fire risk
\end{minipage} & \begin{minipage}[b]{\linewidth}\raggedright
Partially
\end{minipage} & \begin{minipage}[b]{\linewidth}\raggedright
A central target because branch breakers protect wiring, not necessarily
the plug-socket interface.
\end{minipage} \\
\begin{minipage}[b]{\linewidth}\raggedright
Short circuit
\end{minipage} & \begin{minipage}[b]{\linewidth}\raggedright
Immediate high-current damage and fire
\end{minipage} & \begin{minipage}[b]{\linewidth}\raggedright
Yes
\end{minipage} & \begin{minipage}[b]{\linewidth}\raggedright
Reviewed for scope delimitation; not the main target of the smart plug.
\end{minipage} \\
\begin{minipage}[b]{\linewidth}\raggedright
Ground fault
\end{minipage} & \begin{minipage}[b]{\linewidth}\raggedright
Shock and leakage hazard
\end{minipage} & \begin{minipage}[b]{\linewidth}\raggedright
Yes
\end{minipage} & \begin{minipage}[b]{\linewidth}\raggedright
Reviewed to distinguish personnel-protection devices from the proposed
socket-level fire-prevention focus.
\end{minipage} \\
\begin{minipage}[b]{\linewidth}\raggedright
Series arc fault
\end{minipage} & \begin{minipage}[b]{\linewidth}\raggedright
Ignition-capable arcing below ordinary breaker trip levels
\end{minipage} & \begin{minipage}[b]{\linewidth}\raggedright
Not always
\end{minipage} & \begin{minipage}[b]{\linewidth}\raggedright
A central target because its current is load-limited and often
waveform-dependent.
\end{minipage} \\
\begin{minipage}[b]{\linewidth}\raggedright
High-resistance Connection
\end{minipage} & \begin{minipage}[b]{\linewidth}\raggedright
Localized contact heating, glowing, melting, ignition
\end{minipage} & \begin{minipage}[b]{\linewidth}\raggedright
No
\end{minipage} & \begin{minipage}[b]{\linewidth}\raggedright
A central target because the contact can overheat without causing a
classical overcurrent event.
\end{minipage} \\
\begin{minipage}[b]{\linewidth}\raggedright
Overvoltage
\end{minipage} & \begin{minipage}[b]{\linewidth}\raggedright
Equipment damage and stress
\end{minipage} & \begin{minipage}[b]{\linewidth}\raggedright
Partially
\end{minipage} & \begin{minipage}[b]{\linewidth}\raggedright
Included because voltage supervision improves protection and event
interpretation.
\end{minipage} \\
\begin{minipage}[b]{\linewidth}\raggedright
Undervoltage
\end{minipage} & \begin{minipage}[b]{\linewidth}\raggedright
Equipment stress, motor overheating, unreliable operation
\end{minipage} & \begin{minipage}[b]{\linewidth}\raggedright
Partially
\end{minipage} & \begin{minipage}[b]{\linewidth}\raggedright
Included as a supporting protection and measurement-quality condition.
\end{minipage} \\
\midrule\noalign{}
\endhead
\bottomrule\noalign{}
\endlastfoot
\end{longtable}
}

Table 2.1 shows that the proposed smart plug is not intended to replace
breakers, GFCIs, AFCIs, or surge protective devices. Instead, it is
positioned as a complementary protection layer for progressive
fire-inducing faults that may remain dangerous before conventional
protection is forced to operate. Standard circuit breakers are already
effective at mitigating immediate, high-current events; therefore, this
research excludes short circuits and high-current parallel arc faults as
primary targets. The scope instead concentrates on more insidious fault
conditions such as overload behavior that creates slow cumulative heat
damage, series arc faults that operate below ordinary breaker
thresholds, and contact heating caused by degraded plug-socket
interfaces.

It also clarifies why the study uses the term contact heating rather
than treating high-resistance connection as the final hazard itself. A
high-resistance connection is the electrical condition, but the
fire-relevant consequence is localized heating. Contact heating is
therefore defined in this study as the hazardous thermal state that
occurs when a poor or loose localized connection at the plug-socket
interface generates intense heat from ordinary operating current. This
is the specific thermal pathway the prototype aims to detect because it
can develop without immediate conventional protection response.

\section{Related Studies}\label{sec:related-studies}

With the theoretical basis of residential faults and the limitations of
standard protective devices established, this section evaluates existing
research to determine the most appropriate technical direction for the
study. The review begins with non-intrusive load monitoring because load
context affects how waveform disturbances should be interpreted. It then
moves to TinyML feasibility, arc-fault detection features,
threshold-based overload protection, receptacle heating studies, TinyML
training strategy, and validation practices.

\subsection{Non-Intrusive Load Monitoring and Load-Family Identification Studies}\label{non-intrusive-load-monitoring-and-load-family-identification-studies}

Recent non-intrusive load monitoring studies support the idea that
useful appliance information can be extracted from centralized
electrical measurements without installing dedicated sensors on every
device. {[}8{]} shows that NILM can be used not only for energy
disaggregation but also for automation-oriented services in residential
buildings, while the PLAID dataset in {[}40{]} demonstrates that
plug-load identification from sampled voltage and current waveforms is
feasible across many household appliances and operating instances.
Together, these studies justify the use of one electrical measurement
stream for both protection and context recognition {[}8{]}, {[}40{]}.

Feature-based NILM studies are especially relevant because they show
that load identification does not always require large end-to-end
deep-learning models. {[}41{]} uses steady-state current waveform
features in a rule-based framework, while {[}42{]} demonstrates that
percentage THD can serve as a meaningful descriptor of load patterns in
smart-home environments. These findings support classifying appliances
into load families, such as resistive, inductive, rectifier/SMPS,
phase-angle-controlled, and brush-motor loads, before interpreting
abnormal waveform behavior {[}41{]}, {[}42{]}.

The importance of load-family recognition in this study is therefore
protective, not merely descriptive. Healthy household loads can
naturally produce commutation spikes, conduction gaps, or harmonic
distortion that may resemble fault behavior if they are evaluated
without context. The NILM literature supports the claim that family
recognition provides a necessary interpretive layer before a
protection-oriented detector decides whether a waveform disturbance is
intrinsic to the load or evidence of a developing electrical fault
{[}8{]}, {[}40{]}--{[}42{]}.

Since load-family recognition requires embedded interpretation of
waveform features, the next section examines whether machine-learning
inference can realistically be performed on a microcontroller-class
smart plug.

\subsection{TinyML Feasibility on Resource-Constrained Microcontrollers}\label{tinyml-feasibility-on-resource-constrained-microcontrollers}

The feasibility of running machine-learning models on
microcontroller-class hardware is well established in the TinyML
literature. The review in {[}32{]} explains that successful deployment
depends on a hardware-aware workflow in which data preparation, feature
design, model selection, and resource limits are considered together
rather than separately. This is directly relevant to the study, since
the proposed smart plug must perform sensing, signal processing,
inference, actuation, and communication inside a constrained embedded
system {[}32{]}.

This feasibility is reinforced by {[}43{]}, which presents TensorFlow
Lite Micro as an inference framework designed for systems with severe
memory and runtime constraints, and by {[}10{]}, which shows that
TinyML-style anomaly detection can be built around low-cost runtime
operations when the processing pipeline is designed carefully. These
studies support a compact, feature-based embedded pipeline over a large
end-to-end model that would be harder to justify on a small
microcontroller {[}10{]}, {[}43{]}.

The practical implication is that the strongest embedded design is not
necessarily the most complex one, but the one that preserves physical
interpretability and remains deployable under real runtime limits. This
supports the study's hybrid philosophy: deterministic threshold logic
for conditions with clear physical limits, and lightweight inference for
ambiguous waveform cases in which feature interaction and load context
materially improve detection quality {[}10{]}, {[}32{]}, {[}43{]}.

Having established that TinyML deployment is feasible, the review now
focuses on the most difficult fault class for the proposed system:
series arc detection. Unlike overload and socket heating, arc faults
often require multi-feature interpretation because their current
magnitude alone may not be sufficient.

\subsection{Arc-Fault Detection Studies and the Relevance of Multi-Domain Features}\label{arc-fault-detection-studies-and-the-relevance-of-multi-domain-features}

Recent arc-fault literature consistently shows that low-voltage series
arcs are difficult to detect because they are load-dependent,
stochastic, and often masked by healthy nonlinear appliances. The review
in {[}33{]} identifies recurring challenges in AI-based arc-fault
detection, including limited data, false alarms, missed detections, and
poor generalization across loads. This is important to the study because
it shows that arc detection should be framed as a constrained,
multi-feature classification problem rather than as a single-threshold
decision on one waveform quantity {[}33{]}.

Feature-oriented studies point in the same direction. {[}34{]} combines
current-fluctuation features with zero-current features to capture the
concealed and random character of series arcs. {[}35{]} examines THD as
a discriminant for series-arc behavior in single-phase networks, while
{[}36{]} combines high-frequency noise sensing with zero-crossing period
analysis in a multi-load setting. More recent works such as {[}37{]},
{[}38{]}, and {[}39{]} also rely on multi-feature fusion, even though
their final models differ. The common message across these studies is
that useful arc evidence comes from combining complementary indicators
of distortion, randomness, interruption, and context rather than from
relying on current magnitude alone {[}34{]}--{[}39{]}.

Another consistent lesson from this literature is that useful arc
information is not confined to a single frequency band or a single
analysis domain. Some studies exploit high-frequency content, some rely
on zero-current or cycle-randomness descriptors, and others combine
several feature families before classification. For the study, this
supports organizing the current-signal review into two coherent groups:
waveform disturbance and conduction intermittence. These groups are not
independent mechanisms, but complementary views of the same arc process
{[}34{]}--{[}39{]}, {[}49{]}.

\subsubsection{Waveform Disturbance}\label{sec:waveform-disturbance}

Waveform-disturbance features describe how arcing injects broadband
distortion, cycle asymmetry, and unstable spectral content into the
current waveform. A useful starting point is the family of
frequency-band energy features. {[}49{]} performs FFT-based integration
of current content above 100Hz and shows that arc states increase
integrated spectral content, while {[}50{]} uses time-frequency
decomposition and sub-band energy to characterize serial arc behavior.
These studies support the logic of the midband residual ratio, which
isolates spectral energy in a selected midband after suppressing the
dominant fundamental structure, then normalizes that energy against a
baseline current reference so that the resulting ratio is more
comparable across loads {[}49{]}, {[}50{]}.

A second waveform-disturbance descriptor is spectral change over time.
{[}49{]} treats cycle-to-cycle difference and current-cycle similarity
as ways to capture arc randomness, while {[}51{]} includes spectral flux
in a broader waveform-monitoring feature space. For the study, spectral
flux is therefore best understood as a frame-to-frame measure of
spectral instability rather than as a direct measure of arc magnitude.
This makes it useful for identifying the non-stationary spectral
injections associated with unstable conduction {[}49{]}, {[}51{]}.

Waveform disturbance also includes half-cycle shape distortion. {[}49{]}
explicitly notes that, under arc conditions, the current can exhibit
missing waveforms, positive-negative irregularity, and semi-axis
asymmetry. That observation supports the use of half-cycle asymmetry in
the study, where the positive and negative portions of the waveform are
compared as a compact way to quantify unequal conduction between half
cycles {[}49{]}.

For harmonic distortion, the literature supports a contextual rather
than absolute interpretation. {[}35{]} investigated THD for series-arc
detection, while {[}42{]} shows that THD is also informative for
distinguishing healthy load patterns, and {[}52{]} explicitly uses THD
together with the amplitudes of the third and fifth harmonics, RMS, RMS
change rate, and pulse-derived features. This means THD should not be
treated as a standalone proof of arcing; instead, it is better used as
one disturbance feature inside a broader decision process that already
accounts for the natural nonlinearity of certain appliance families
{[}35{]}, {[}42{]}, {[}52{]}.

Finally, high-frequency energy elevation is one of the most common
spectral signatures of series arcing. {[}53{]} designed a band-pass
filter specifically to isolate arc-related content and then developed an
algorithm using arc-signal energy and pulse counts, while {[}36{]} also
relies on high-frequency noise as part of its detection logic. These
studies support the study's use of a high-frequency energy-change
feature that tracks whether the upper-band energy share increases
relative to a baseline or broader band, thereby capturing the broadband
elevation that often accompanies arcing {[}36{]}, {[}53{]}.

Waveform-disturbance features are useful because they capture the
spectral and harmonic effects of arcing. However, arcing is also a
conduction problem: the current may extinguish, restrike, and become
discontinuous.

\subsubsection{Conduction Intermittence}\label{sec:conduction-intermittence}

Conduction-intermittence features describe how series arcs repeatedly
extinguish and restrike, producing irregular current continuity even
when the overall load current remains modest. {[}34{]} is the clearest
support for this line of reasoning because it explicitly computes a
zero-current time-proportional coefficient and combines it with a
current-fluctuation measure. That makes {[}34{]} the proper anchor for
the study's low-current runtime and low-current ratio features: both are
intended to quantify how much of the observation window is spent in weak
or interrupted conduction {[}34{]}.

The same physical logic supports dispersion-based zero-cross metrics. In
a healthy waveform, zero-cross intervals remain relatively regular,
whereas in an arc state, they become more erratic because extinction and
restrike disturb the timing of conduction around the zero region.
{[}36{]} relies directly on zero-crossing period analysis, and {[}34{]}
likewise centers its method on zero-current behavior. Hence, the study
used the term zero-cross variance, which is derived from the dispersion
of zero-cross intervals as outlined by studies {[}34{]} and {[}36{]}.

Changes in RMS from one analysis interval to the next also belong to the
conduction-intermittence family when they are interpreted as evidence of
unstable conduction rather than simple load switching. {[}52{]}
explicitly includes RMS and the change rate of RMS between adjacent
cycles in its lightweight arc-fault method. This provides direct support
for the study's use of an absolute delta-of-RMS feature to capture
cycle-to-cycle instability that is too persistent and irregular to be
explained by healthy steady-state conduction {[}52{]}.

{[}54{]} introduces a moving Z-score method calculated from dual FFT
bands. Building on this concept, the current study employs an absolute
z-deviation from a rolling RMS baseline, adapting the normalization
principle outlined in {[}54{]}. This feature generates a load-relative
deviation score that helps maintain meaningful thresholds across
appliances with significantly different baseline currents {[}54{]}.

The arc-fault literature supports feature-based classification because
series arcs are ambiguous and load-dependent. In contrast, overload is
more directly tied to current magnitude and heating duration.

\subsection{Studies on Overload Monitoring and Threshold-Based Protection}\label{studies-on-overload-monitoring-and-threshold-based-protection}

Threshold-based domestic protection remains relevant because overload is
a physical stress problem before it is a machine-learning problem.
{[}44{]} presents a smart overload trip system that measures current and
disconnects supply when the current exceeds a preset value, while
{[}45{]} presents an ESP32-based smart multiplug with current and
temperature monitoring, remote control, and protective thresholding.
These studies support the study's use of direct and interpretable
current-limit logic for overload protection {[}44{]}, {[}45{]}.

Their importance to the study is architectural. They show that the
smart-plug form factor can already accommodate sensing, communication,
and protective switching. The present study therefore extends an already
credible embedded architecture, but directs it toward a broader
protection role by combining threshold protection with context-aware
fault interpretation rather than stopping at current monitoring alone
{[}44{]}, {[}45{]}.

Since overload protection addresses current-related thermal stress, the
next related fault concern is localized socket heating. Unlike overload,
contact heating may occur even when total load current remains within a
range that does not immediately force upstream breaker operation.

\subsection{Receptacle Heating Studies and Early Socket-Temperature Trip}\label{receptacle-heating-studies-and-early-socket-temperature-trip}

The most important heating studies for this project are the ones showing
that poor contact is a local and progressive fire mechanism. {[}24{]}
demonstrated that loose terminal conditions can produce severe
overheating and receptacle damage across many test configurations.
{[}17{]} showed that receptacles and attachment plugs subjected to
glowing conditions could suffer severe material damage, while the
current path itself could continue functioning. These findings strongly
support the idea that a connection can remain electrically functional
even when thermally dangerous.

{[}46{]} simulated poor electric contact between plug and receptacle and
showed that contact pressure strongly influences temperature rise and
fire risk. {[}47{]} experimentally studied the ignition process caused
by poor electrical contact and described a staged mechanism progressing
from resistive heating and pyrolysis to arc-accelerated ignition.
{[}7{]} also investigated glowing contact in electrical fires and
highlighted that glowing phenomena may persist for hours while
transferring hazardous heat to surrounding materials. Taken together,
these studies show that socket heating is not a secondary symptom but a
primary ignition pathway.

The temperature-threshold question is best interpreted through studies
of hidden versus measured temperature. {[}30{]} argued that there can be
a large difference between overheated contact-spot temperature and
external surface temperature. {[}31{]} again demonstrated that external
package temperature can under-read internal contact temperature. This is
the strongest scientific support for an early trip threshold such as
40°C at an external NTC sensor. The argument is not that all outlets
melt at 40°C, but that a 40°C external reading can already correspond to
a hidden contact region operating far above that value. Considering the
literature on standards, degradation, thermal lag, and indirect sensing,
along with the limitations of the NTC sensor, the 40°C threshold is
treated in this study as a conservative estimated tripping temperature.

The related studies on overload and socket heating justify deterministic
threshold protection, while the arc-fault studies justify feature-based
inference. The next section connects these two approaches by discussing
how TinyML training should be structured for embedded fault detection.

\subsection{TinyML Training}\label{tinyml-training}

The reviewed literature also supports the overall training strategy used
by the study. First, feature nomination should begin from fault physics
rather than from blind feature accumulation. NILM and arc studies
consistently favor features that describe harmonic content, waveform
stability, conduction gaps, restrikes, and cross-cycle irregularity
because these are closer to the physical differences between healthy
nonlinear loads and unstable electrical faults {[}34{]}--{[}37{]},
{[}41{]}, {[}42{]}. This supports the study's decision to treat waveform
disturbance and conduction intermittence as two coherent feature
families.

Second, the training data should represent the operating regimes the
detector will actually see. That means not only arc windows, but also
normal steady-state windows, turn-on transients, and load-family
diversity. The literature on load dependence and false positives makes
this especially important: a detector trained only on clean resistive
examples will not generalize well to motors, rectifiers, or phase-angle
loads {[}8{]}, {[}33{]}, {[}40{]}. Thus, the training pipeline used in
the project, which captures start, steady, and arc segments, follows a
structure already justified by the literature.

Third, model selection in TinyML is not solely about raw test accuracy.
The reviews by {[}32{]} and {[}43{]} both imply that model choice must
account for runtime footprint, deterministic behavior, and deployment
simplicity. For this reason, a practical workflow is to screen plausible
feature subsets using a relatively fast ensemble method, then finalize a
compact classifier such as a Random Forest or other lightweight model
for deployment. This kind of staged search is consistent with the
broader TinyML literature because it balances exploratory breadth with
embedded realism.

Finally, the load-family context should be trained separately or jointly
in a way that informs the final arc decision. The literature strongly
suggests that context should influence thresholds, persistence rules, or
posterior interpretation rather than being treated as an isolated
auxiliary label {[}8{]}, {[}33{]}, {[}41{]}. In practical terms, this
means a smart plug should recognize whether it is observing a resistive
heater, a motor, a rectifier, or a phase-angle-controlled load before
deciding how aggressively to interpret impulsive waveform disturbances.

After establishing how the model should be trained, the next concern is
how the system should be validated. Because the proposed device
addresses different fault classes, the validation procedure must also be
divided according to the physical behavior of each fault.

\subsection{Testing and Validation}\label{testing-and-validation}

The validation literature supports using controlled, repeatable
fault-induction methods rather than waiting for naturally occurring
household faults. {[}48{]} constructs a dedicated AC arc-fault platform
with a fixed electrode, a moving electrode driven by a stepper motor,
and multiple household loads, which is the kind of repeatable setup
needed for arc-data collection and algorithm evaluation. This supports
the use of controlled arc-generation procedures in the study {[}48{]}.

Validation should also extend beyond one appliance type. The
load-diversity argument in NILM and arc-fault work implies that good
performance on one heater or one motor is insufficient. A
protection-oriented smart plug should therefore be validated across
several load families, nuisance events, and non-fault transients so that
its specificity is tested alongside its sensitivity {[}33{]}, {[}40{]},
{[}41{]}, {[}48{]}.

For overload and contact-heating validation, the literature supports a
layered approach. Threshold-based smart-plug studies validate whether
trip logic responds consistently when safety limits are exceeded, while
poor-contact studies justify separate thermal validation for localized
heating risk. The resulting implication for the study is that overload,
arc detection, and contact heating should not be merged into one generic
accuracy test; they require distinct procedures, reference measurements,
and nuisance checks {[}44{]}--{[}48{]}.

Table 2.2 summarizes how the reviewed literature maps to the feature and
decision families used in the study. It groups the quantities into
waveform disturbance, conduction intermittence, and deterministic stress
variables to clarify why the prototype combines feature-based TinyML
inference with threshold-based protection logic.

\subparagraph{\texorpdfstring{\textbf{Table 2.2.} Feature Families
Relevant in the
Study}{Table 2.2. Feature Families Relevant in the Study}}\label{table-2.2.-feature-families-relevant-in-the-study}

{\def\LTcaptype{none} % do not increment counter
\begin{longtable}[]{@{}
  >{\raggedright\arraybackslash}p{(\linewidth - 4\tabcolsep) * \real{0.2024}}
  >{\raggedright\arraybackslash}p{(\linewidth - 4\tabcolsep) * \real{0.3228}}
  >{\raggedright\arraybackslash}p{(\linewidth - 4\tabcolsep) * \real{0.4015}}@{}}
\toprule\noalign{}
\begin{minipage}[b]{\linewidth}\centering
\textbf{Feature Family}
\end{minipage} & \begin{minipage}[b]{\linewidth}\centering
\textbf{Representative Quantities}
\end{minipage} & \begin{minipage}[b]{\linewidth}\centering
\textbf{Literature}
\end{minipage} \\
\begin{minipage}[b]{\linewidth}\raggedright
Waveform disturbance
\end{minipage} & \begin{minipage}[b]{\linewidth}\raggedright
Midband Residual Ratio;

Spectral Flux at Mid to High Frequency;

Halfcycle Assymetry;

Current Total Harmonic Distortion;

High Frequency Energy-Change
\end{minipage} & \begin{minipage}[b]{\linewidth}\raggedright
Arc studies repeatedly report that series arcs inject random harmonic
and high-frequency disturbance, while NILM studies show that healthy
loads also have family-specific spectral signatures, making these
features informative but context-dependent {[}34{]}-{[}39{]}, {[}41{]},
{[}42{]}, {[}49{]}-{[}53{]}.
\end{minipage} \\
\begin{minipage}[b]{\linewidth}\raggedright
Conduction intermittence
\end{minipage} & \begin{minipage}[b]{\linewidth}\raggedright
Low-Current Runtime and Low Current Ratio;

Zero-Cross Variance;

Absolute Delta of Current RMS;

Absolute Z-Deviation from Baseline
\end{minipage} & \begin{minipage}[b]{\linewidth}\raggedright
These quantities represent restrikes, current interruption, partial
conduction, and unstable cycle timing, which are physically consistent
with series arcs and weak-contact behavior {[}34{]}, {[}36{]}, {[}37{]},
{[}52{]}, {[}54{]}.
\end{minipage} \\
\begin{minipage}[b]{\linewidth}\raggedright
Deterministic stress variables
\end{minipage} & \begin{minipage}[b]{\linewidth}\raggedright
Irms; temperature estimate; voltage sag or abnormal supply condition
\end{minipage} & \begin{minipage}[b]{\linewidth}\raggedright
These variables directly encode thermal or electrical stress and
therefore remain essential even when machine learning is added. They
improve interpretability and provide hard safety overrides for overload
and heating events {[}3{]}, {[}27{]}-{[}29{]}, {[}44{]}, {[}45{]}.
\end{minipage} \\
\midrule\noalign{}
\endhead
\bottomrule\noalign{}
\endlastfoot
\end{longtable}
}

Table 2.2 shows that the proposed feature set is not arbitrary. The
waveform-disturbance features represent spectral and harmonic effects of
arcing, while the conduction-intermittence features represent the
extinguish-and-restrike behavior of unstable current flow. The
deterministic stress variables remain necessary because some faults,
particularly overload and socket heating, are better handled through
direct physical thresholds than through classification alone. This table
therefore supports the hybrid architecture of the study: machine
learning is used where waveform interpretation is ambiguous, while
deterministic logic is retained where physical safety limits are clear.

\section{Synthesis}\label{synthesis}

The reviewed literature supports the proposed smart plug most strongly
when the device is framed as a complementary socket-level protection
layer. Existing smart-plug and overload-monitoring studies already show
that compact embedded devices can measure current and temperature,
communicate with users, and actuate protective switching {[}44{]},
{[}45{]}. At the same time, poor-contact fire studies show that the
plug-receptacle interface is itself a meaningful origin of heating and
ignition risk {[}46{]}, {[}47{]}. Taken together, these findings justify
a protection device that monitors hazards locally at the socket rather
than relying only on upstream protection behavior {[}44{]}--{[}47{]}.

The literature also shows that overload, contact heating, and series
arcing are related but not equivalent fault classes. Overload is
fundamentally a current-stress problem and is therefore well-suited to
deterministic threshold logic. Contact heating is a localized thermal
problem caused by degraded or poor electrical contact and therefore
requires direct local temperature supervision. Series arcing, by
contrast, is load-dependent and stochastic; it often appears as a
mixture of spectral disturbance, zero-current behavior, and cycle
irregularity rather than as a dramatic increase in current magnitude.
This is why the most defensible design is a hybrid one: hard thresholds
for direct electrical or thermal stress, and feature-based inference for
ambiguous waveform cases {[}34{]}--{[}39{]}, {[}44{]}--{[}47{]}.

A second synthesis point concerns context. NILM and load-identification
studies establish that healthy appliances have distinct waveform
signatures, while arc-fault studies show that many healthy nonlinear
behaviors can resemble fault signatures if load context is ignored.
Therefore, load-family recognition is not merely an added analytics
feature; it is part of the protection logic itself because it helps
separate load-intrinsic nonlinearity from fault-induced non-stationarity
{[}8{]}, {[}33{]}--{[}42{]}.

A third synthesis point concerns implementation feasibility. The TinyML
literature shows that on-device inference is realistic on
microcontroller-class hardware when the pipeline is built around compact
features, controlled memory use, and deployment-aware model selection.
That makes the proposed architecture technically plausible: a smart plug
can combine deterministic rules with a lightweight model, provided the
feature computation, inference path, and safety logic are kept small,
interpretable, and robust {[}10{]}, {[}32{]}, {[}43{]}.

The gap addressed by the study is therefore clear. Prior work has
separately shown the feasibility of smart-plug monitoring, the need for
local supervision of poor-contact heating, the value of load-aware
waveform interpretation, and the practicality of TinyML on constrained
hardware. What remains comparatively underdeveloped is a single embedded
device that unifies these ideas into one protection-oriented platform.
The proposed TinyML Smart Plug addresses that gap by combining
socket-level thermal supervision, threshold-based overload protection,
load-family recognition, and multi-feature series-arc interpretation in
one embedded system {[}8{]}, {[}10{]}, {[}32{]}--{[}47{]}.
